(1) 令天空角坐标 \(\boldsymbol\alpha\)，圆盘半径 \(a_\theta=\delta/2\)，亮度 \(I(\boldsymbol\alpha)=I_0\) 在 \(|\boldsymbol\alpha|\le a_\theta\) 内，其余为零。归一化复相干度为 \(\gamma(\boldsymbol B)=\int I(\boldsymbol\alpha)e^{-i(2\pi/\lambda)\boldsymbol B\cdot\boldsymbol\alpha}d^2\alpha/\int I\,d^2\alpha\)。取基线方向为极轴后，角积分变为 Bessel 函数 \(J_0\)，并且轴对称源使 \(\gamma\) 为实偶函数。

(2) 径向积分 \(\int_0^{a_\theta}\rho J_0(k\rho)d\rho=a_\theta J_1(ka_\theta)/k\)，其中 \(k=2\pi B/\lambda\)。由圆盘面积归一得到 \(\gamma(B)=2J_1(\pi B\delta/\lambda)/(\pi B\delta/\lambda)\)。实验条纹可见度是 \(V(B)=|\gamma(B)|\)，即会失去 \(J_1\) 在不同瓣间的正负相位。圆盘各方向等价，因此仅依赖基线长度而不依赖其天球平面方向。

(3) \(J_1(x)\) 第一正零点 \(x_1\simeq3.831706\)，故第一个相干零点满足 \(\pi B_0\delta/\lambda=x_1\)，即 \(B_0\simeq1.21967\,\lambda/\delta\)。测得基线零点并知道波长可反求恒星角直径 \(\delta\simeq1.22\lambda/B_0\)，假设关键是均匀圆盘亮度。对边缘昏暗的真实恒星，零点系数与亮度律有关，不能未经模型检验直接套用。

(4) 小 \(x=\pi B\delta/\lambda\) 时 \(J_1(x)=x/2-x^3/16+O(x^5)\)，因而 \(\gamma=1-x^2/8+O(x^4)\)，初始可见度从1按 \(B^2\) 下降。只测 \(|\gamma|\) 对中心对称源无法恢复丢失的符号相位，甚至对一般非对称源会出现相位恢复歧义；需闭合相位或其他附加相位参考，不能从单条基线的条纹对比度推断天体在天空中的绝对方位。
